\documentclass[12pt]{article}
\usepackage[a4paper,margin=1in]{geometry}
\usepackage{amsmath,amssymb}
\usepackage{enumitem}
\usepackage{hyperref}

\begin{document}

\section*{\textbf{CSE400: Fundamentals of Probability in Computing}}

\subsection*{\textbf{Lecture Scribe 5: Bayes’ Theorem, Random Variables, and Probability Mass Function (PMF)}}

\textbf{Instructor:} Dhaval Patel, PhD \\
\textbf{Date:} January 20, 2026

\hrule
\vspace{1em}

\subsection*{\textbf{1. Outline}}

This lecture covers the following topics:

\begin{enumerate}
    \item Bayes’ Theorem
    \begin{itemize}
        \item Weighted Average of Conditional Probabilities
        \item Learning by Example
        \item Formal Introduction: Law of Total Probability and Bayes Theorem
    \end{itemize}

    \item Random Variables
    \begin{itemize}
        \item Motivation and Concept
        \item Examples
    \end{itemize}

    \item Probability Mass Function (PMF)
    \begin{itemize}
        \item Concept and Examples
    \end{itemize}

    \item Class Participation Quiz
\end{enumerate}

\hrule
\vspace{1em}

\section*{\textbf{2. Bayes’ Theorem}}

\subsection*{\textbf{2.1 Weighted Average of Conditional Probabilities}}

Let $A$ and $B$ be events. We may express $A$ as:

\[
A = AB \cup AB^c
\]

This is because for an outcome to be in $A$, it must either be in both $A$ and $B$, or be in $A$ but not in $B$.

Since $AB$ and $AB^c$ are mutually exclusive, by \textbf{Axiom 3}, we have:

\[
Pr(A) = Pr(AB) + Pr(AB^c)
\]

Using conditional probability form:

\[
Pr(A) = Pr(A|B)Pr(B) + Pr(A|B^c)[1 - Pr(B)]
\]

\subsubsection*{\textbf{Key Result}}

The probability of event $A$ is a \textbf{weighted average} of conditional probabilities, where the weights are given by the probability of the event on which it is conditioned.

\hrule
\vspace{1em}

\section*{\textbf{3. Bayes’ Theorem: Learning by Example}}

\subsection*{\textbf{3.1 Example 3.1 (Part 1/2)}}

An insurance company believes that people can be divided into two classes:

\begin{itemize}
    \item Accident prone
    \item Not accident prone
\end{itemize}

The statistics show:

\begin{itemize}
    \item An accident-prone person will have an accident within 1 year with probability $0.4$
    \item A not accident-prone person will have an accident within 1 year with probability $0.2$
\end{itemize}

Assume $30\%$ of the population is accident prone.

\subsubsection*{\textbf{Question}}

What is the probability that a new policyholder will have an accident within a year of purchasing a policy?

\subsubsection*{\textbf{Solution}}

Let:

\begin{itemize}
    \item $A_1$ denote the event that the policyholder will have an accident within a year
    \item $A$ denote the event that the policyholder is accident prone
\end{itemize}

Then:

\[
Pr(A_1) = Pr(A_1|A)Pr(A) + Pr(A_1|A^c)Pr(A^c)
\]

Substituting values:

\[
Pr(A_1) = (0.4)(0.3) + (0.2)(0.7)
\]

\[
= 0.26
\]

\subsubsection*{\textbf{Final Answer}}

\[
Pr(A_1) = 0.26
\]

\subsection*{\textbf{3.2 Example 3.1 (Part 2/2)}}

Suppose that a new policyholder has an accident within a year of purchasing a policy.

\subsubsection*{\textbf{Question}}

What is the probability that he or she is accident prone?

\subsubsection*{\textbf{Solution}}

The desired probability is:

\[
Pr(A|A_1) = \frac{Pr(AA_1)}{Pr(A_1)}
\]

Using:

\[
Pr(AA_1) = Pr(A)Pr(A_1|A)
\]

Thus:

\[
Pr(A|A_1) = \frac{Pr(A)Pr(A_1|A)}{Pr(A_1)}
\]

Substituting:

\[
Pr(A|A_1) = \frac{(0.3)(0.4)}{0.26}
\]

\[
= \frac{6}{13}
\]

\subsubsection*{\textbf{Final Answer}}

\[
Pr(A|A_1) = \frac{6}{13}
\]

\hrule
\vspace{1em}

\section*{\textbf{4. Formal Introduction: Law of Total Probability and Bayes Formula}}

\subsection*{\textbf{4.1 Law of Total Probability (Formula 3.4)}}

This is known as the \textbf{law of total probability}.

\subsection*{\textbf{4.2 Bayes Formula (Proposition 3.1)}}

Using:

\[
Pr(AB_i) = Pr(B_i|A)Pr(A)
\]

we get the Bayes Formula.

This is known as the \textbf{Bayes Formula (Proposition 3.1)}.

\subsubsection*{\textbf{Interpretation}}

\begin{itemize}
    \item $Pr(B_i)$ is the \textbf{apriori probability} (probabilities formed from self-evident or presupposed models)
    \item $Pr(B_i|A)$ is the \textbf{posteriori probability} (probabilities derived or calculated after observing certain events) of event $B_i$ given $A$
\end{itemize}

\hrule
\vspace{1em}

\section*{\textbf{5. Bayes Formula: Learning by Example}}

\subsection*{\textbf{5.1 Example 3.2}}

Suppose we have 3 cards that are identical in form, except:

\begin{itemize}
    \item Both sides of the first card are red
    \item Both sides of the second card are black
    \item One side of the third card is red and the other side is black
\end{itemize}

The 3 cards are mixed up in a hat, and 1 card is randomly selected and put down on the ground.

If the upper side of the chosen card is colored red, what is the probability that the other side is colored black?

\subsubsection*{\textbf{Solution}}

Let:

\begin{itemize}
    \item $RR$: chosen card is all red
    \item $BB$: chosen card is all black
    \item $RB$: chosen card is red-black
    \item $R$: upturned side of chosen card is red
\end{itemize}

We need:

\[
Pr(RB|R)
\]

Using conditional probability:

\[
Pr(RB|R) = \frac{Pr(RB \cap R)}{Pr(R)}
\]

\[
= \frac{Pr(R|RB)Pr(RB)}{Pr(R|RR)Pr(RR) + Pr(R|RB)Pr(RB) + Pr(R|BB)Pr(BB)}
\]

Substitute probabilities:

\[
= \frac{(1/2)(1/3)}{(1)(1/3) + (1/2)(1/3) + (0)(1/3)}
\]

\[
= \frac{1}{3}
\]

\subsubsection*{\textbf{Final Answer}}

\[
Pr(RB|R) = \frac{1}{3}
\]

\subsubsection*{\textbf{Incorrect Reasoning Mentioned}}

Many incorrectly guess $1/2$ by assuming that the two possibilities (all-red and red-black) are equally likely. This is incorrect because the outcomes are not equally likely.

If each card is treated as two distinct sides, then there are 6 equally likely outcomes:

\[
R_1, R_2, B_1, B_2, R_3, B_3
\]

Since the other side is black only if the outcome is $R_3$, the desired probability equals the conditional probability of $R_3$ given that either $R_1$, $R_2$, or $R_3$ occurred, which equals:

\[
\frac{1}{3}
\]

\hrule
\vspace{1em}

\section*{\textbf{6. Random Variables}}

\subsection*{\textbf{6.1 Motivation and Concept}}

When an experiment is performed, we are frequently interested mainly in some function of the outcome rather than the actual outcome itself.

\subsubsection*{\textbf{Experiment 1}}

In dice tossing, the focus is often on the sum of the two dice rather than the individual values.

\subsubsection*{\textbf{Experiment 2}}

When flipping a coin, the interest may lie in the total number of heads without focusing on the exact head-tail sequence.

These real-valued functions defined on the sample space are known as \textbf{Random Variables}.

\begin{itemize}
    \item Values are determined by the outcomes of an experiment.
    \item Probabilities are assigned to possible values of random variables.
\end{itemize}

\subsection*{\textbf{6.2 Visualization of Distribution}}

The distribution of a random variable can be visualized as a bar diagram:

\begin{itemize}
    \item The x-axis represents the values the random variable can take on.
    \item The height of the bar at value $a$ is the probability:
\end{itemize}

\[
Pr[X = a]
\]

Each probability can be computed by looking at the probability of the corresponding event in the sample space.

\hrule
\vspace{1em}

\section*{\textbf{7. Random Variables: Examples}}

\subsection*{\textbf{7.1 Example 1 (Three Fair Coin Tosses)}}

Suppose the experiment consists of tossing 3 fair coins.

Let $Y$ denote the number of heads that appear.

Then $Y$ is a random variable taking values:

\[
0, 1, 2, 3
\]

with probabilities:

\[
P\{Y = 0\} = P\{(t,t,t)\} = \frac{1}{8}
\]

\[
P\{Y = 1\} = P\{(t,t,h),(t,h,t),(h,t,t)\} = \frac{3}{8}
\]

\[
P\{Y = 2\} = P\{(t,h,h),(h,t,h),(h,h,t)\} = \frac{3}{8}
\]

\[
P\{Y = 3\} = P\{(h,h,h)\} = \frac{1}{8}
\]

Since $Y$ must take one of the values $0$ through $3$, we must have:

\[
\sum_y P(Y=y) = 1
\]

\hrule
\vspace{1em}

\section*{\textbf{8. Probability Mass Function (PMF)}}

\subsection*{\textbf{8.1 Discrete Random Variable}}

A random variable that can take on at most a countable number of possible values is said to be \textbf{discrete}.

\subsection*{\textbf{8.2 Definition of Probability Mass Function}}

Let $X$ be a discrete random variable with range:

\[
R_X = x_1, x_2, x_3, \dots
\]

(finite or countably infinite).

The function:

\[
p(x) = Pr(X=x)
\]

is called the \textbf{Probability Mass Function (PMF)} of $X$.

Since $X$ must take one of the values $x_k$, we have:

\[
\sum_k p(x_k) = 1
\]

\hrule
\vspace{1em}

\section*{\textbf{9. PMF Examples}}

\subsection*{\textbf{9.1 Example: Two Independent Tosses of a Fair Coin}}

(Example is included as shown in lecture material.)

\subsection*{\textbf{9.2 Example: PMF Given in Exponential Form}}

The probability mass function of a random variable $X$ is given by:

\[
p(i) = c\lambda^i,\quad i = 0,1,2,\dots
\]

where $\lambda$ is some positive value.

\subsubsection*{\textbf{Question}}

Find:

\begin{enumerate}
    \item $P\{X=0\}$
    \item $P\{X>2\}$
\end{enumerate}

(As per lecture material, the example is included exactly as written.)

\hrule
\vspace{1em}

\section*{\textbf{10. Class Participation}}

Students are instructed to switch to \textbf{Campuswire} for class participation and discussion.

\hrule
\vspace{1em}

\section*{\textbf{End of Lecture Scribe 5}}

\textbf{Topic Coverage Completed:} \\
Bayes’ Theorem, Law of Total Probability, Bayes Formula, Random Variables, PMF.

\end{document}
